\begin{ouverture}[spectre]{16}{Lumière et spectres}{Ondes et signaux}{Observer · Modéliser · Expérimenter · Expliquer}
\objectif{Utiliser la propagation rectiligne et la vitesse de la lumière.}
\objectif{Distinguer spectre continu et spectre de raies.}
\objectif{Caractériser une lumière monochromatique par sa longueur d’onde.}
\prerequis{Proportionnalité, conversions et lecture de graphiques.}
\end{ouverture}
\essentiel
\essSection{1}{Des rayons pour modéliser}
Dans un milieu transparent et homogène, la lumière se propage en ligne droite. Un rayon représente une direction et un sens de propagation, pas un objet matériel. Pour voir un objet, de la lumière provenant de cet objet doit entrer dans l'œil.
\begin{definition}[Vitesse]Dans le vide, la vitesse de la lumière vaut exactement \qty{299792458}{m.s^{-1}} ; pour les exercices du livre dans le vide ou l'air, $c\simeq\qty{3,00e8}{m.s^{-1}}$. La durée de parcours vaut $\Delta t=d/c$.\end{definition}
La lumière parcourt environ un mètre en quelques nanosecondes ; elle est beaucoup plus rapide que le son dans l'air.
\essSection{2}{Décomposer une lumière}
Un prisme ou un réseau sépare les radiations d'une lumière polychromatique : on obtient un spectre. Une radiation monochromatique est caractérisée par sa longueur d'onde dans le vide ou l'air, souvent exprimée en nanomètres. Le visible s'étend approximativement de 400 à 800 nm selon les conventions et la sensibilité de l'œil.
\schema[corrige]{spectres}
\begin{propriete}[Corps chaud]Un corps chaud dense émet un spectre continu. Quand sa température augmente, le spectre s'enrichit vers les courtes longueurs d'onde du visible ; la couleur apparente évolue. Une source gazeuse excitée peut présenter un spectre de raies caractéristique des espèces présentes.\end{propriete}
\essSection{3}{Identifier par des raies}
Pour exploiter un spectre, comparer les positions des raies à celles de références avec la même échelle de longueurs d'onde. Une seule raie commune ne suffit pas toujours à identifier une espèce ; rechercher un ensemble cohérent.
\begin{savoirfaire}[Méthode]Lire les unités, relever plusieurs longueurs d'onde, comparer aux valeurs de référence, puis formuler une conclusion avec la précision de lecture. Un spectre schématique sans axe gradué ne permet pas une lecture numérique précise.\end{savoirfaire}
Une lumière colorée peut être monochromatique ou résulter de plusieurs radiations ; la couleur apparente seule ne donne pas la composition spectrale.
\essSection{4}{Expérimenter et modéliser}
\begin{experience}[Observer un spectre]Observer avec un spectroscope une lampe adaptée puis une source à raies sous encadrement. Comparer les spectres, repérer continu ou discontinu, relever l'échelle. Ne jamais regarder directement le Soleil ni un faisceau laser. Pour une mesure numérique, vérifier l'étalonnage en longueur d'onde.\end{experience}
\exercices
\rubrique{Prendre confiance}
\begin{blocLivre}{ex-16-1}\begin{exercice}[1]{Une valeur à connaître}
On utilise la vitesse de la lumière dans l’air pour un calcul scolaire.
\begin{questions}
\item Donner sa valeur approchée en SI.
\item La comparer qualitativement à 340 m/s.
\end{questions}
\end{exercice}\end{blocLivre}
\begin{blocLivre}{ex-16-2}\begin{exercice}[1]{Des nanomètres aux mètres}
Une radiation a pour longueur d’onde 550 nm.
\begin{questions}
\item Exprimer cette valeur en mètres.
\item Le nanomètre désigne-t-il une durée ?
\end{questions}
\end{exercice}\end{blocLivre}
\begin{blocLivre}{ex-16-3}\begin{exercice}[1]{Continuité du spectre}
La figure compare une bande continue et un ensemble de raies.
\schema{spectres}
\begin{questions}
\item Associer les deux types de spectres.
\item Lequel peut correspondre à un corps chaud dense ?
\end{questions}
\end{exercice}\end{blocLivre}
\begin{blocLivre}{ex-16-4}\begin{exercice}[1]{Un trajet lumineux}
Une impulsion lumineuse parcourt 900 m dans l’air ; $c=3,00\times10^8$ m/s.
\begin{questions}
\item Calculer sa durée de propagation.
\item Exprimer la durée en microsecondes.
\end{questions}
\end{exercice}\end{blocLivre}
\rubrique{S’entraîner}
\begin{blocLivre}{ex-16-5}\begin{exercice}[2]{Reconnaître des raies}
La source X possède des raies à 450, 520 et 650 nm. Les références sont A : 450, 520, 650 nm ; B : 430, 550, 680 nm.
\begin{questions}
\item Quelle référence correspond à X ?
\item Pourquoi comparer plusieurs raies ?
\end{questions}
\end{exercice}\end{blocLivre}
\begin{blocLivre}{ex-16-6}\begin{exercice}[2]{Chauffer davantage}
Un filament chaud passe d’une lueur rouge à une lumière plus blanche.
\begin{questions}
\item Comment évolue qualitativement son spectre visible ?
\item Faut-il un spectre de raies pour expliquer cette évolution ?
\end{questions}
\end{exercice}\end{blocLivre}
\begin{blocLivre}{ex-16-7}\begin{exercice}[2]{Monochromatique ou non}
Une source verte possède deux raies à 530 et 550 nm.
\begin{questions}
\item Est-elle monochromatique ?
\item La couleur verte suffit-elle à donner une seule longueur d’onde ?
\end{questions}
\end{exercice}\end{blocLivre}
\begin{blocLivre}{ex-16-8}\begin{exercice}[2]{Voir un objet}
Un livre est éclairé dans une salle.
\begin{questions}
\item Quel trajet de lumière permet sa vision ?
\item L’œil doit-il envoyer un rayon vers le livre ?
\end{questions}
\end{exercice}\end{blocLivre}
\rubrique{Raisonner et approfondir}
\begin{blocLivre}{ex-16-9}\begin{exercice}[3]{Une distance astronomique}
La lumière d’une source met 500 s pour nous parvenir ; $c=3,00\times10^8$ m/s.
\begin{questions}
\item Calculer la distance.
\item Exprimer cette distance en kilomètres.
\end{questions}
\end{exercice}\end{blocLivre}
\begin{blocLivre}{ex-16-10}\begin{exercice}[3]{Un aller-retour lumineux}
Un télémètre reçoit une impulsion réfléchie après 40,0 ns. On adopte $c=3,00\times10^8$ m/s.
\begin{questions}
\item Calculer la distance à la cible.
\item Pourquoi diviser par deux ?
\end{questions}
\end{exercice}\end{blocLivre}
\begin{blocLivre}{ex-16-11}\begin{exercice}[3]{La lecture d’un mélange}
Un spectre contient toutes les raies de A et deux raies supplémentaires absentes de A.
\begin{questions}
\item Quelle présence est compatible avec le spectre ?
\item Peut-on affirmer que la source contient seulement A ?
\end{questions}
\end{exercice}\end{blocLivre}
\begin{blocLivre}{ex-16-12}\begin{exercice}[3]{Éclair et tonnerre}
Un éclair est vu puis le tonnerre est entendu 3,0 s plus tard. $v_s=340$ m/s.
\begin{questions}
\item Estimer la distance en négligeant le temps lumineux.
\item Justifier cette approximation à partir de l’ordre de grandeur.
\end{questions}
\end{exercice}\end{blocLivre}

